% Texte de l'oral de soutenance MémoTDSI — compilation : pdflatex texte-oral.tex (deux fois)
\documentclass[11pt,a4paper]{article}

\usepackage[utf8]{inputenc}
\usepackage[T1]{fontenc}
\usepackage[french]{babel}
\usepackage{lmodern}
\usepackage[a4paper,top=2.6cm,bottom=2.4cm,left=2.3cm,right=2.3cm,headheight=28pt]{geometry}
\usepackage{graphicx}
\usepackage[table]{xcolor}
\usepackage{tabularx}
\usepackage{booktabs}
\usepackage{longtable}
\usepackage{array}
\usepackage{enumitem}
\usepackage{fancyhdr}
\usepackage{lastpage}
\usepackage[most]{tcolorbox}
\usepackage[hidelinks]{hyperref}

% ------------------------------------------------------------------ couleurs de l'UCAD
\definecolor{ucadbleu}{HTML}{2993CD}
\definecolor{ucadnuit}{HTML}{0B4F8A}
\definecolor{ucador}{HTML}{C9A227}
\definecolor{doux}{HTML}{5A6475}
\definecolor{fondbleu}{HTML}{EEF5FB}

% ------------------------------------------------------------------ en-tête et pied de page
\pagestyle{fancy}
\fancyhf{}
\fancyhead[L]{\raisebox{-0.35\height}{\includegraphics[height=22pt]{logo-ucad.png}}\hspace{6pt}%
  {\small\color{ucadnuit}\textbf{Université Cheikh Anta Diop de Dakar}}}
\fancyhead[R]{\small\color{doux}Soutenance MémoTDSI — texte de l'oral}
\fancyfoot[L]{\footnotesize\color{doux}Master Mathématiques, Cryptographie et Sécurité — TDSI}
\fancyfoot[R]{\footnotesize\color{doux}\thepage\ / \pageref{LastPage}}
\renewcommand{\headrule}{{\color{ucadbleu}\hrule height 1.2pt}\vspace{1pt}{\color{ucador}\hrule height 0.4pt}}
\renewcommand{\footrulewidth}{0pt}

% ------------------------------------------------------------------ mise en forme
\setlength{\parindent}{0pt}
\setlength{\parskip}{4pt}

% Titre d'une partie d'orateur
\newcommand{\orateur}[3]{%
  \clearpage
  \begin{tcolorbox}[colback=ucadnuit,colframe=ucadnuit,coltext=white,arc=2pt,
      left=10pt,right=10pt,top=8pt,bottom=8pt]
    {\large\textbf{#1}}\hfill{\color{ucador}\textbf{Diapos #2}}\\[2pt]
    {\small #3}
  \end{tcolorbox}\vspace{4pt}}

% Une diapo : numéro, titre, texte à dire (entre guillemets)
\newcommand{\diapo}[3]{%
  \begin{tcolorbox}[enhanced,breakable,colback=white,colframe=ucadbleu!45,
      borderline west={3pt}{0pt}{ucadbleu},arc=2pt,
      left=8pt,right=8pt,top=5pt,bottom=6pt]
    {\color{ucadnuit}\textbf{Diapo #1} — #2}\par\vspace{3pt}
    \og #3\fg{}
  \end{tcolorbox}\vspace{2pt}}

% Phrase de passage au suivant
\newcommand{\passage}[1]{%
  \vspace{2pt}{\color{doux}\textit{Passage :} \og #1\fg{}}}

\newcolumntype{Y}{>{\raggedright\arraybackslash}X}

\begin{document}

% ================================================================== page de titre
\thispagestyle{empty}
\begin{center}
  {\small\textbf{RÉPUBLIQUE DU SÉNÉGAL}}\\[1pt]
  {\footnotesize\itshape Un Peuple — Un But — Une Foi}\\[14pt]
  \includegraphics[width=3.4cm]{logo-ucad.png}\\[10pt]
  {\LARGE\bfseries\color{ucadnuit} UNIVERSITÉ CHEIKH ANTA DIOP DE DAKAR}\\[6pt]
  {\large\bfseries Faculté des Sciences et Techniques}\\[3pt]
  {\normalsize Département de Mathématiques et Informatique}\\[3pt]
  {\normalsize Master Mathématiques, Cryptographie et Sécurité — TDSI}\\[10pt]
  {\color{ucadbleu}\rule{0.9\linewidth}{1.6pt}}\\[-6pt]
  {\color{ucador}\rule{0.9\linewidth}{0.5pt}}\\[40pt]

  {\Huge\bfseries\color{ucadnuit} MémoTDSI}\\[10pt]
  {\Large La mémothèque de la filière TDSI}\\[6pt]
  {\large\color{doux} Texte de l'oral de soutenance — 38 diapositives, 6 orateurs}\\[50pt]

  \begin{tcolorbox}[colback=fondbleu,colframe=ucadbleu!50,arc=2pt,width=0.82\linewidth,
      left=12pt,right=12pt,top=8pt,bottom=8pt]
    \centering
    \textbf{Équipe}\\[2pt]Mouhamadou Siradji Sow \\ Ahmadou Ba \\ Alseny DIALLO \\ Aldemba Wade \\ Amy Ndiaye \\ Mamadou Aminata Gueye\\[6pt]
    \textbf{Année universitaire :} 2026–2027\\[4pt]
    \textbf{Durée :} environ 25 minutes de présentation et 8 minutes de démonstration
  \end{tcolorbox}

  \vfill
  {\footnotesize\color{doux} Université Cheikh Anta Diop de Dakar — BP 5005 Dakar-Fann, Sénégal — www.ucad.sn}
\end{center}

% ================================================================== organisation
\clearpage
\section*{\color{ucadnuit}Organisation}

Les six orateurs se partagent les 38 diapos. Chacun termine par une phrase de passage qui annonce le suivant.
Le texte est fait pour être dit à voix haute : on peut le reformuler avec ses mots, mais on garde les
chiffres exacts (225 tests, 113 classes, 600\,000 itérations, 80 bits).

\vspace{6pt}
\rowcolors{2}{fondbleu}{white}
\begin{tabularx}{\linewidth}{@{}l l Y r@{}}
  \toprule
  \textbf{Orateur} & \textbf{Diapos} & \textbf{Thème} & \textbf{Durée} \\
  \midrule
  M1 — Mouhamadou Siradji Sow & 1 à 6   & Présentation, plan, circuit, espace étudiant & 4 min \\
  M2 — Ahmadou Ba & 7 à 11  & Enseignant, planification, PV, convocations  & 4 min \\
  M3 — Alseny DIALLO & 12 à 18 & Nouveautés, pile technique, architecture     & 4 min \\
  M4 — Aldemba Wade & 19 à 25 & Le code Java : service, objets, JDBC, JavaFX & 5 min \\
  M5 — Amy Ndiaye & 26 à 30 & Cryptographie et protection des données      & 4 min \\
  M6 — Mamadou Aminata Gueye & 31 à 38 & QR codes, TOTP, PDF, tests, démonstration    & 4 min + démo \\
  \bottomrule
\end{tabularx}

% ================================================================== M1
\orateur{M1 — Mouhamadou Siradji Sow}{1 à 6}{L'application : présentation, plan, acteurs, circuit, portail, espace étudiant}

\diapo{1}{MémoTDSI}{Bonjour. Nous sommes l'équipe M1 à M6 du Master Mathématiques, Cryptographie et
Sécurité. Nous présentons MémoTDSI, la mémothèque de la filière TDSI. C'est une application de bureau qui
suit un mémoire de bout en bout, du dépôt par l'étudiant jusqu'à sa publication avec un procès-verbal signé
électroniquement. Comme nous sommes en master sécurité, la sécurité a été pensée dès le départ, pas ajoutée
à la fin.}

\diapo{2}{Plan}{Notre présentation suit quatre parties. D'abord l'application, avec Ahmadou Ba et moi-même.
Ensuite les nouveautés pour la scolarité et l'architecture, avec Alseny DIALLO. Puis le code Java, avec
Aldemba Wade, et la cryptographie, avec Amy Ndiaye. Enfin, Mamadou Aminata Gueye présentera les documents
vérifiables et la double authentification, puis lancera la démonstration.}

\diapo{3}{Trois acteurs, un circuit}{Trois acteurs : l'étudiant, l'enseignant et la scolarité. Comme dans le
vrai système d'information d'une université, chacun ne fait que ce qui relève de son rôle. L'étudiant
dépose, l'encadreur donne son avis, et la scolarité organise le reste : la campagne, les soutenances et les
procès-verbaux.}

\diapo{4}{Le circuit d'un mémoire}{Quatre étapes. Un : l'étudiant dépose son PDF, qui est contrôlé, chiffré,
et dont on calcule l'empreinte SHA-256. Deux : l'encadreur donne un avis favorable ou demande des
corrections. Trois : la scolarité planifie la soutenance. Quatre : après la soutenance, la note donne la
mention et le PV est signé. En cas de corrections, ou si la note est inférieure à 10, on revient à
l'étudiant. L'étudiant ne fixe ni la date ni la mention.}

\diapo{5}{Un portail universitaire}{L'interface ressemble à un portail d'université : photo du campus à la
connexion, menu avec des icônes, cloche de notifications et textes en grand. Le tableau de bord dit à chacun
ce qu'il doit faire maintenant, par exemple \og 2 mémoires à évaluer\fg{}. Pour la démonstration, les
comptes de test se remplissent en un clic.}

\diapo{6}{Espace étudiant}{L'étudiant suit son mémoire étape par étape : une pastille de couleur montre
l'état, avec la date de soutenance et le résultat. Il imprime son attestation de dépôt, qui porte
l'empreinte SHA-256 de son PDF, et sa convocation officielle.}

\passage{Je laisse la parole à Ahmadou Ba, qui va vous montrer le côté enseignant et scolarité.}

% ================================================================== M2
\orateur{M2 — Ahmadou Ba}{7 à 11}{Enseignants et scolarité : planification, procès-verbal, convocations, pilotage}

\diapo{7}{Espace enseignant}{L'encadreur ne note pas. Il autorise la soutenance, ou demande des corrections
avec un motif obligatoire. Avant toute décision, l'application vérifie que le PDF n'a pas été modifié. Dans
\og Mes jurys\fg{}, l'enseignant voit son rôle dans chaque soutenance (président, examinateur ou encadreur)
et imprime sa convocation.}

\diapo{8}{La scolarité planifie}{La scolarité choisit la date, la salle et le jury. Avant de valider, le
service vérifie les règles : un horaire entre 7\,h et 20\,h, au moins un examinateur, et l'encadreur ne peut
pas présider. Il vérifie aussi qu'aucune salle et qu'aucun enseignant ne sont pris deux fois au même moment.
L'étudiant et chaque membre du jury sont ensuite notifiés.}

\diapo{9}{Le procès-verbal}{Après la soutenance, la scolarité saisit la note et la mention s'affiche en
direct, selon le barème LMD. Pour signer le PV, il faut retaper son mot de passe. Avec 10 ou plus, le
mémoire est publié ; en dessous, il est ajourné et revient à l'étudiant. La diffusion peut être libre,
restreinte, ou confidentielle avec un embargo.}

\diapo{10}{Les convocations}{Pour chaque soutenance, l'application rédige une lettre officielle par
destinataire : l'étudiant, le président, l'examinateur et l'encadreur. Elle porte l'en-tête de l'UCAD, le
tableau de la soutenance et la place pour le cachet. La scolarité imprime toutes les lettres ; chaque autre
personne ne peut obtenir que la sienne.}

\diapo{11}{Pilotage}{La scolarité voit tous les mémoires d'un coup d'œil : des compteurs par étape, des
filtres et des pastilles de couleur.}

\passage{Alseny DIALLO va maintenant vous présenter deux nouveautés et la façon dont le code est organisé.}

% ================================================================== M3
\orateur{M3 — Alseny DIALLO}{12 à 18}{Nouveautés et architecture : campagne par niveau, import Excel, pile technique}

\diapo{12}{Partie 2 (intercalaire)}{Voici deux nouveautés demandées par la scolarité.}

\diapo{13}{Campagne de dépôt par niveau}{Chaque niveau, de L1 à M2, a maintenant sa propre campagne. La
scolarité l'ouvre ou la ferme, fixe une date limite, et la prolonge ou l'écourte avec les boutons $+7$ et
$-7$ jours. Les étudiants du niveau sont prévenus à chaque changement, par exemple : \og date limite du
niveau M2 prolongée du 30/11 au 07/12\fg{}. Dans le code, une \texttt{EnumMap<Niveau, Campagne>} contient
la campagne de chaque niveau. Au passage, nous avons corrigé un bug : le sélecteur de date de JavaFX
n'enregistrait une date tapée au clavier qu'après un appui sur Entrée.}

\diapo{14}{Comptes et import Excel}{Il y a trois types de comptes : étudiant, enseignant et personnel
administratif, chacun avec son téléphone. La scolarité peut importer un fichier Excel entier ; chaque ligne
donne un résultat : créé, réinscrit, mis à jour ou erreur. Le point intéressant : un fichier \texttt{.xlsx}
n'est qu'une archive ZIP de fichiers XML. Nous le lisons avec le JDK, sans Apache POI, et nous nous
protégeons contre la bombe ZIP et contre l'attaque XXE.}

\diapo{15}{Partie 3 (intercalaire)}{Passons sous le capot : comment l'application est construite en Java.}

\diapo{16}{La pile technique}{Java 21, JavaFX 21 pour l'interface, SQLite avec JDBC pour les données, Maven
pour construire et JUnit 5 pour tester. Seulement 4 dépendances : tout le reste est écrit par nous avec le
JDK, y compris la cryptographie, la lecture Excel, les QR codes et les PDF. Le projet compte 113 classes,
environ 15\,500 lignes et 225 tests.}

\diapo{17}{Organisation du code}{Un paquet par couche : \texttt{controller} pour un contrôleur par écran,
\texttt{service} pour les règles métier, \texttt{dao} pour l'accès à la base, \texttt{model} pour les objets
et \texttt{util} pour les outils communs. Une couche n'appelle que celle du dessous ; un écran ne touche
donc jamais la base.}

\diapo{18}{Architecture}{C'est le schéma MVC, complété par une couche Service et une couche DAO. Le point le
plus important est en rouge : les droits sont vérifiés dans les services. Cacher un bouton ne suffit pas ;
si quelqu'un appelait la méthode autrement, le service refuserait quand même.}

\passage{Aldemba Wade va vous montrer le code Java lui-même, à partir d'un exemple réel.}

% ================================================================== M4
\orateur{M4 — Aldemba Wade}{19 à 25}{Le code Java : parcours d'un clic, service, objets, DAO, JavaFX, threads, exceptions}

\diapo{19}{Le parcours d'un clic}{Prenons un exemple concret : la scolarité planifie une soutenance. Le clic
sur OK arrive au contrôleur, qui appelle le service. Le service vérifie le droit, puis les règles. Le DAO
cherche les conflits en SQL, puis enregistre tout dans une seule transaction, et les notifications partent.
Si une règle est violée, une exception remonte : le message s'affiche, la fenêtre reste ouverte et rien
n'est enregistré.}

\diapo{20}{Le code du service}{Voici le vrai code de \texttt{planifier}. La première ligne est
\texttt{exigerAdmin()}, le contrôle du rôle. Chaque règle lève ensuite une exception avec un message clair.
Le statut du mémoire est un enum : seul un mémoire \texttt{AUTORISE} peut être programmé. Comme le service
ne dépend pas de l'écran, nos tests l'appellent directement, avec chaque rôle.}

\diapo{21}{Modèle objet et Java moderne}{Les états sont des enums, comme \texttt{StatutMemoire} : le
compilateur vérifie les valeurs, ce qui évite les fautes de frappe dans des chaînes. Les valeurs immuables
sont des records, comme \texttt{Campagne}. Nous utilisons aussi les nouveautés de Java : switch expressions,
pattern matching avec \texttt{instanceof}, et text blocks pour écrire le SQL.}

\diapo{22}{Accès aux données : DAO + JDBC}{\texttt{GenericDAO<T>} est une interface générique qui regroupe les
opérations communes. Le point à retenir, c'est la transaction : \texttt{setAutoCommit(false)}, puis
\texttt{commit}, ou \texttt{rollback} en cas d'erreur. C'est tout ou rien. Toutes les requêtes passent par
\texttt{PreparedStatement}, donc aucune injection SQL n'est possible. Le schéma de la base évolue par
migrations numérotées, de V2 à V10.}

\diapo{23}{L'interface JavaFX}{Le fichier FXML décrit l'écran. Le contrôleur reçoit automatiquement les
composants grâce à \texttt{@FXML} : un \texttt{fx:id} dans le FXML correspond à un champ dans le contrôleur.
Le style tient dans une feuille CSS à jetons de couleur, et le thème sombre ne redéfinit que ces jetons.}

\diapo{24}{Ne jamais figer l'interface}{JavaFX n'a qu'un seul fil pour l'affichage. Un calcul long, comme un
PBKDF2 à 600\,000 itérations, figerait la fenêtre. Notre classe \texttt{TacheFond} fait donc ce calcul sur un
autre fil, puis revient sur le fil JavaFX pour afficher le résultat. Pendant le calcul, les boutons sont
grisés, ce qui évite un double clic.}

\diapo{25}{Les erreurs, côté Java}{Chaque exception a un sens. \texttt{IllegalArgumentException} : une saisie
invalide. \texttt{IllegalStateException} : une étape du circuit non respectée. \texttt{SecurityException} :
un droit refusé. \texttt{SQLException} : une erreur technique. Dans les fenêtres, un multi-catch affiche le
message, et \texttt{ev.consume()} garde la fenêtre ouverte.}

\passage{Amy Ndiaye va maintenant vous expliquer comment la cryptographie protège les données.}

% ================================================================== M5
\orateur{M5 — Amy Ndiaye}{26 à 30}{La cryptographie : JDK, protection des données, signature, falsification, journal}

\diapo{26}{La cryptographie avec le JDK}{Aucune bibliothèque externe : l'API standard du JDK suffit. PBKDF2
protège les mots de passe, AES-GCM chiffre les fichiers, HMAC chaîne le journal et Ed25519 signe les
documents. Les clés sont rangées dans un coffre PKCS12, séparé de la base.}

\diapo{27}{Protection des données}{Le PDF déposé est chiffré en AES-256-GCM. La clé est dans un coffre séparé,
et l'empreinte SHA-256 est vérifiée à chaque ouverture. Le mot de passe est haché avec PBKDF2, 600\,000
itérations et un sel de 16 octets.}

\diapo{28}{Signature du procès-verbal}{Le PV est signé avec Ed25519. La signature couvre l'empreinte du PDF,
la note, la mention, la date et le signataire. La clé privée de la scolarité est chiffrée par son mot de
passe : pour signer, il faut le retaper.}

\diapo{29}{Falsifier un résultat ? Détecté}{C'est notre démonstration préférée. Si l'on change la mention
directement dans la base, par exemple \og Bien\fg{} en \og Très bien\fg{}, la fiche affiche aussitôt
\og signature invalide\fg{}.}

\diapo{30}{Un journal infalsifiable}{Chaque action sensible est journalisée. Chaque ligne est liée à la
précédente par un HMAC. Si l'on modifie, insère ou supprime une ligne, la vérification de la chaîne le
détecte.}

\passage{Mamadou Aminata Gueye va présenter nos dernières fonctions de confiance, puis lancer la démonstration.}

% ================================================================== M6
\orateur{M6 — Mamadou Aminata Gueye}{31 à 38}{Confiance et démonstration : documents vérifiables, QR, TOTP, PDF, tests}

\diapo{31}{Partie 4 (intercalaire)}{Dernière partie : des documents qu'on peut vérifier, et une connexion
renforcée.}

\diapo{32}{Documents vérifiables}{Une convocation imprimée se falsifie facilement. Chaque convocation, PV ou
attestation porte donc un QR code qui contient ses informations et leur signature Ed25519, ainsi qu'un code
de 16 caractères imprimé à côté. L'écran \og Vérifier un document\fg{} est accessible même sans compte. Il
répond : authentique, plus valable (par exemple si la soutenance a été déplacée), falsifié ou inconnu.}

\diapo{33}{Un encodeur QR écrit en Java}{Nous avons écrit l'encodeur QR nous-mêmes, d'après la norme
ISO~18004. Il protège les données par un code correcteur de Reed-Solomon dans le corps fini GF($2^8$), ce
qui est directement notre cours. Il est validé sur 67 QR codes, identiques module par module à une
implémentation de référence.}

\diapo{34}{Double authentification TOTP}{Après le mot de passe, l'application demande un code à 6 chiffres,
comme Google Authenticator, qui change toutes les 30 secondes. Ce code est un HMAC-SHA1 du secret et de
l'heure. Un code ne sert qu'une fois, et cinq codes faux verrouillent le compte. Nous l'avons testé avec les
vecteurs officiels de la RFC~6238.}

\diapo{35}{Export PDF sans bibliothèque}{Les aperçus ont un bouton \og Enregistrer en PDF\fg{}. Nous écrivons
le fichier PDF nous-mêmes, objet par objet, sans iText ni PDFBox, et un test vérifie que sa structure est
correcte.}

\diapo{36}{Stratégie de tests}{225 tests JUnit 5, chacun sur une base SQLite jetable. Ils couvrent le circuit
complet, les droits de chaque rôle, la falsification, les vecteurs officiels des normes et l'import Excel.
GitHub Actions les relance à chaque push.}

\diapo{37}{Qualité}{En résumé : 225 tests, aucune bibliothèque de cryptographie externe, trois rôles
strictement séparés, et un thème clair ou sombre.}

\diapo{38}{Merci}{Merci de votre attention. Place à la démonstration.}

\vspace{6pt}
\begin{tcolorbox}[colback=fondbleu,colframe=ucador,arc=2pt,left=10pt,right=10pt,top=6pt,bottom=6pt,
    title={\textbf{Démonstration (environ 8 minutes)}},coltitle=ucadnuit,colbacktitle=fondbleu,
    fonttitle=\normalsize]
\begin{enumerate}[leftmargin=*,itemsep=2pt]
  \item Babacar NDIAYE (\texttt{babs@mcs} / \texttt{babs123}) : sa convocation avec QR code, puis
        \og Enregistrer en PDF\fg{}.
  \item M. SAMB (\texttt{samb@mcs} / \texttt{samb123}) : donne son avis sur un mémoire.
  \item La scolarité (\texttt{admin@mcs} / \texttt{admin123}) : $+7$ jours sur la campagne M2, planification
        d'une soutenance, saisie du PV signé.
  \item \og Vérifier un document\fg{} : coller le texte du QR, résultat authentique ; changer une lettre du
        nom, résultat falsifié.
  \item \og Mon profil\fg{} : activer la double authentification avec un téléphone, se déconnecter, puis se
        reconnecter avec le code.
  \item Modifier la mention dans la base SQLite : la fiche affiche \og signature invalide\fg{}.
\end{enumerate}
\end{tcolorbox}

% ================================================================== questions
\clearpage
\section*{\color{ucadnuit}Questions probables du professeur de Java}

Chaque question a un répondant désigné : celui qui a présenté la diapo concernée.

\vspace{4pt}
\rowcolors{2}{fondbleu}{white}
\begin{longtable}{@{}>{\raggedright\arraybackslash}p{0.28\linewidth}>{\raggedright\arraybackslash}p{0.47\linewidth}>{\raggedright\arraybackslash}p{0.19\linewidth}@{}}
  \toprule
  \textbf{Question} & \textbf{Réponse courte} & \textbf{Répond} \\
  \midrule
  \endhead
  Pourquoi pas Spring ou Hibernate ? & Pour une application de bureau sur un seul poste, JDBC et des couches
    écrites à la main suffisent, et montrent qu'on comprend ce que ferait un framework. & Alseny DIALLO \\
  À quoi sert le rollback ? & Sans lui, on pourrait enregistrer une soutenance sans son jury. Avec, c'est
    tout ou rien. & Aldemba Wade \\
  Pourquoi un enum plutôt qu'une chaîne ? & Le compilateur vérifie les valeurs ; une faute de frappe devient
    une erreur de compilation, pas un bug. & Aldemba Wade \\
  Que se passe-t-il si on calcule sur le fil JavaFX ? & La fenêtre se fige pendant le calcul. C'est pour ça
    que \texttt{TacheFond} utilise un autre fil. & Aldemba Wade \\
  Comment empêchez-vous l'injection SQL ? & \texttt{PreparedStatement} partout, et un tri choisi dans une
    liste fermée (un enum). & Aldemba Wade \\
  Où sont vérifiés les droits ? & Dans les services, pas dans les écrans. Les tests appellent les services
    avec un rôle non autorisé et attendent une \texttt{SecurityException}. & Alseny DIALLO \\
  Pourquoi AES-GCM et pas AES-CBC ? & GCM chiffre et authentifie : toute modification du fichier fait échouer
    le déchiffrement. & Amy Ndiaye \\
  Pourquoi 600\,000 itérations de PBKDF2 ? & C'est la recommandation actuelle de l'OWASP pour
    PBKDF2-HMAC-SHA256 : chaque essai coûte cher à un attaquant. & Amy Ndiaye \\
  Comment savez-vous que votre QR code est juste ? & 67 codes comparés module par module à une
    implémentation de référence, plus trois matrices de référence dans les tests. & Mamadou Aminata Gueye \\
  Comment empêcher de réutiliser un code TOTP ? & On mémorise le dernier pas de 30 secondes accepté ; un code
    d'un pas égal ou plus ancien est refusé. & Mamadou Aminata Gueye \\
  \bottomrule
\end{longtable}

\end{document}
